\begin{center}
{\heiti\zihao{3} 基于随机森林与共表达网络融合的啤酒风味预测}\\[0.25em]
{\heiti\zihao{3} 及关键化合物识别}
\end{center}

\vspace{0.6em}

\section{摘要}

啤酒风味是数百种化合物在酿造与陈酿过程中共同作用的结果，具有维度高、样本少、共线性强的特点。本文以 \textbf{250} 个啤酒样本的 \textbf{231} 个化学特征和 \textbf{50} 个感官维度为对象，建立“数据口径统一—\textbf{随机森林}回归—分级标定—\textbf{共表达网络}融合排序”的建模框架，并对模型进行多角度检验。

\textbf{针对问题一}，将感官评分预测写成 \textbf{50} 维多输出回归，采用 \textbf{随机森林}（由多棵决策树平均的集成模型）拟合化合物与感官维度之间的非线性关系，再用\textbf{锚定线性映射与四舍五入}（ALR）把连续预测转成 \textbf{0--6} 级整数。两阶段预注册选型得到随机森林与 \textbf{226} 维特征集 F0，折外 \textbf{R\textsuperscript{2}} 为 \textbf{0.2565}、五折均值为 \textbf{0.2305}，高于同口径基线 \textbf{0.243}。针对 \textbf{0} 级与 \textbf{6} 级为空的现象，本文证明这是\textbf{刻度不可达命题}：折外预测为条件均值，预测与标签的标准差之比为 \textbf{0.444}，方差比约 \textbf{0.197}，与 \textbf{R\textsuperscript{2}} 同阶；锚点冻结时 \textbf{6} 级要求预测越过标签约 \textbf{92} 百分位，故两端等级不可达是口径的数学必然，而非实现缺陷。辅助的\textbf{可读口径}表在保持单调同序（逆序对 \textbf{0}）的前提下使 \textbf{7} 级全部出现。

\textbf{针对问题二}，以 \textbf{226} 个候选化合物为节点、以 \textbf{Spearman} 秩相关绝对值不低于 \textbf{0.30} 为边，构建含 \textbf{3758} 条边的共表达网络；将直接相关得分与邻居加权得分按 \textbf{0.70} 的图项权重融合，为每个感官维度给出 \textbf{Top-3} 关键化合物。\textbf{Bootstrap} 重采样 \textbf{200} 次显示，第一名化合物的平均入选频率为 \textbf{0.6734}，第三名降至 \textbf{0.3611}；\textbf{Jaccard} 重合度为 \textbf{0.378}，说明图项确实改变了排序。以来源品类核对，一致 \textbf{48} 对、不一致 \textbf{86} 对、不可比 \textbf{16} 对，一致率 \textbf{0.3582}，高于独立基线 \textbf{0.2047}。\textbf{梯度提升机}（GBM）全局对照中 \textbf{60/60} 的重合由候选池包含关系恒真，本文只把平均分位 \textbf{0.6201} 对 \textbf{0.4971} 作为有信息量的证据。

\textbf{针对问题三}，题面要求以风格为条件生成化合物谱与感官画像并做跨风格迁移。本文只给出条件生成、化学约束与回评的拟采用路线，不报告任何结果，具体见第 \textbf{6} 章。

最后，本文用缺失扰动、锚点扰动和外部证据检验模型：缺失比例增至 \textbf{20\%} 时折均 \textbf{R\textsuperscript{2}} 由 \textbf{0.2305} 降至 \textbf{0.2121}，分级一致率仍为 \textbf{0.8882}；外部证据 E1、E2、E3 的方向一致率分别为 \textbf{0.9968}、\textbf{0.8098}、\textbf{0.9706}。模型在可解释性与稳健性之间取得了可复核的平衡。

\noindent\textbf{关键词：}啤酒风味预测；随机森林；共表达网络；分级标定；稳定性分析；关键化合物识别；小样本高维建模
